\chapter{Introduction}
\label{chp:introduction}

Container terminals connect maritime and hinterland transport networks. Higher
container volumes and larger vessels require terminal operators to make many
tactical and operational decisions under uncertainty. These decisions concern
truck dispatching, yard-block allocation, resource use, congestion, equipment
failures, and temporary infrastructure outages. Experiments in an operating
terminal are often costly, disruptive, or not feasible. Simulation offers a way
to assess decisions without interrupting terminal operations.

\section{Motivation}
Simulation is used for decision support in container terminal planning and
operations research. It can assess infrastructure investments, operational
policies, scheduling approaches, truck appointment systems, and terminal
redesigns before implementation. Traditional approaches mostly use manually
developed Discrete Event Simulation (DES) models, agent-based models, or
digital twins. Building and maintaining these models requires domain knowledge,
manual work, and calibration \cite{YUN1999DESSimulation,CTSteenken2004,OptimizingCTs2023,Bett2024}.

Process Mining research has studied how simulation models can be built
from event logs. Earlier work showed that control-flow behaviour, activity
durations, arrival patterns, resource allocations, and routing probabilities
can be discovered from process data and used in executable simulation models
\cite{PMRozinat2009}. \textit{Data-aware process simulation} extends this idea.
It uses context attributes, case information, process states, resource
characteristics, and decision logic
\cite{DaSLeno2019,mannhardt2023stochasticProcessesModelling,Vinci2023Inf.Activ.Prob,lopez2024}.

Machine-learning and deep-learning methods can learn complicated behaviour,
but their results are often difficult to inspect, modify, and validate. Hybrid
and white-box approaches combine interpretable models with data-driven
behaviour discovery \cite{runtimeIntegration2024,Vinci2026prosit}. ProSiT uses
probabilistic white-box predictors to build simulation models that can be
inspected and configured for what-if analysis \cite{Vinci2026prosit}.

Evidence on automatically discovered simulation models in logistics systems is
still limited. Existing evaluations mainly concern business processes in
healthcare, financial services, and administration
\cite{Hospital2026Vinci,Vinci2023Inf.Activ.Prob}. Container terminal truck
operations involve physical resource constraints, congestion,
infrastructure dependencies, spatial interactions, and changing process
contexts. It is not yet clear how well a discovered data-aware simulation model
can reproduce these conditions and support decisions. This thesis studies its
fidelity, structural applicability, and limitations using event logs from real
container terminal truck operations.

\section{Problem Statement}
Container terminal operators make decisions about capacity, truck handling,
yard use, policies, and infrastructure changes. Experiments in a real terminal
are often not feasible or involve operational risk, so these decisions are
often studied with simulation. Most container-terminal simulation models are
built manually. Their construction and maintenance require domain knowledge
and modelling work. They may include subjective assumptions about process
behaviour, resource interactions, and routing logic
\cite{CTSteenken2004,OptimizingCTs2023}.

Data-aware process simulation discovers simulation models from historical
event logs. This may reduce modelling work and keep the model closer to the
recorded process \cite{Camargo2020,lopez2024,Vinci2023Inf.Activ.Prob}. Before
such a model can support decisions, it must reproduce observed operations,
remain configurable, and produce stable results. Earlier studies have shown
that data-aware simulation can be applied to business processes
\cite{Vinci2023Inf.Activ.Prob,Hospital2026Vinci}. Its use in logistics systems
with congestion, competing resources, and context-dependent behaviour is less
clear. This thesis evaluates the model's fidelity, use, and limitations for
container-terminal truck operations.

\section{Research Objective} The research objective of this thesis is to
evaluate the fidelity, structural applicability, and limitations of an
automatically discovered data-aware process simulation model for container
terminal truck operations. Trustworthiness motivates the study, but one
historical case study cannot establish it as a general property. The thesis
aims to:

\begin{itemize} 
\item quantify how well a discovered ProSiT model reproduces unseen CTB truck-handling behaviour;
\item evaluate whether the deliberately restrictive training trace language covers held-out cases while respecting the physical sequence of a truck visit;
\item determine how increasing the contextual expressiveness of the simulation configuration affects held-out fidelity; and
\item examine how the model responds to transparent operational interventions and what those responses reveal about the limits of what-if decision support.
\end{itemize} 

The study uses event data from container terminal truck operations for process
discovery and simulation model generation.

\subsection{Research Contributions} The contributions of this thesis are:

\begin{enumerate} 
\item \textbf{Domain evaluation}\\ The thesis provides one of the first
evaluations of data-aware process simulation for container terminal truck
operations.
\item \textbf{Method}\\ It identifies challenges that arise when data-aware
simulation is transferred from business processes to logistics systems with
physical constraints and congestion.
\item \textbf{Model comparison}\\ It measures how different levels of context
affect held-out simulation fidelity and variation across seeds.
\item \textbf{Decision-support limits}\\ It uses operational interventions to
identify what the discovered model can and cannot represent in what-if
analysis.
\end{enumerate} 

\section{Research Questions} To achieve the research objective, the thesis addresses the following research questions:

\begin{enumerate} 
\item \textbf{RQ1 (Fidelity):} To what extent can a ProSiT simulation model discovered from historical CTB truck-operation data reproduce the observed behaviour of an unseen temporal hold-out?
\item \textbf{RQ2 (Structural applicability):} To what extent can the discovered process structure represent held-out CTB truck-handling behaviour while respecting the physical sequential constraints of a truck visit?
\item \textbf{RQ3 (Contextual expressiveness):} How does increasing the contextual expressiveness of the simulation model affect held-out simulation fidelity in the CTB truck process?
\item \textbf{RQ4 (What-if decision support):} How does the discovered simulation model respond to transparent operational interventions, and what do these responses reveal about its suitability and limitations for what-if decision support?
\end{enumerate} 

\section{Thesis Outline} 
\label{sec:thesis_outline} 

Chapter~\ref{chp:background} introduces container-terminal operations, Process
Mining, Business Process Simulation, data-aware and white-box simulation, and
methods for simulation validation.

Chapter~\ref{chp:case_study_and_data} describes HHLA Container Terminal
Burchardkai, the data sources, and the pipeline that turns raw records into the
simulation event log.

Chapter~\ref{chp:experiments_and_framework} defines the train/test split,
process and simulation discovery, Monte-Carlo replications, fidelity measures,
two what-if scenarios, the transition-duration analysis, and the required
ProSiT workarounds.

Chapter~\ref{chap:results} reports the event-log characteristics, Petri-net
conformance, arrival process, held-out fidelity of the ProSiT configurations,
per-activity results, discovered resource capacities, transition durations, and
what-if results. Confidence intervals or stated uncertainty accompany the
empirical results.

Chapter~\ref{chp:conclusion} answers the research questions, discusses the
limitations, and gives directions for future work.

